% ============================================================================
% Fichier  : annexes/Ann_C_modeles.tex
% Titre    : Modèles de Documents
% Version  : 1.0 -- novembre 2026
% ============================================================================

\chapter{Modèles de Documents}
\label{ann:modeles}

Les modèles suivants reprennent les documents utilisés dans le cours. Ils se
recopient tels quels dans un tableur ou un traitement de texte, et
s'adaptent à l'organisation.

\section*{M1. Registre des actions}

\begin{tableaucours}{L{1.6cm}L{1.6cm}L{2.0cm}L{3.6cm}L{3.0cm}}{%
    \tableauHeader{N°} &
    \tableauHeader{Date, heure} &
    \tableauHeader{Auteur} &
    \tableauHeader{Action} &
    \tableauHeader{Résultat, empreinte}
}
001 & & & & \\
002 & & & & \\
\end{tableaucours}

\noindent Chapitres~\ref{chap:prologue} et~\ref{chap:enquete}.

\section*{M2. Formulaire de chaîne de possession}

\begin{enumerate}
    \item Identifiant de la pièce et numéro de scellé.
    \item Description~: nature, marque, modèle, numéro de série, capacité.
    \item Provenance~: système, emplacement, propriétaire.
    \item Recueil~: date, heure, lieu, personne, méthode, outil et version.
    \item Empreintes~: algorithme, valeur, heure, témoin.
    \item Conditionnement~: enveloppe, sac, scellé, photographie.
    \item Transferts~: remettant, destinataire, date, heure, motif, état du
          scellé, signatures.
    \item Examens~: personne, date, opération, copie utilisée.
    \item Conservation~: lieu, conditions.
    \item Sort final~: restitution, destruction, archivage, autorisation.
\end{enumerate}

\noindent Chapitre~\ref{chap:preuve}.

\section*{M3. Lettre de mission d'enquête}

\begin{itemize}
    \item Mandant et enquêteur~; date~; fondement (contrat, réquisition).
    \item Questions de l'enquête, numérotées.
    \item Périmètre~: systèmes, comptes, période.
    \item Autorisations~: signataires, hébergeurs et prestataires concernés.
    \item Règles de conservation des pièces et de remise à l'autorité
          judiciaire.
    \item Livrables et destinataires.
\end{itemize}

\noindent Chapitres~\ref{chap:investigation} et~\ref{chap:normes-droit}.

\section*{M4. Matrice des hypothèses concurrentes}

\begin{tableaucours}{L{5.6cm}cccc}{%
    \tableauHeader{Pièce} &
    \tableauHeader{H1} & \tableauHeader{H2} & \tableauHeader{H3} &
    \tableauHeader{H4}
}
& & & & \\
\end{tableaucours}

\noindent C~: compatible~; I~: incompatible~; N~: neutre. On retient
l'hypothèse qui a le moins d'incompatibilités (Chapitre~\ref{chap:investigation}).

\section*{M5. Tableau des horloges}

\begin{tableaucours}{L{3.6cm}cccL{2.6cm}}{%
    \tableauHeader{Source} &
    \tableauHeader{Fuseau} &
    \tableauHeader{Dérive} &
    \tableauHeader{Précision} &
    \tableauHeader{Synchronisée}
}
& & & & \\
\end{tableaucours}

\noindent Chapitre~\ref{chap:chronologie}.

\section*{M6. Fiche de constat de vulnérabilité}

\begin{itemize}
    \item Identifiant (CVE), faiblesse (CWE), produit et version.
    \item Actif~: nom, propriétaire, exposition, importance.
    \item Preuve du constat~: méthode (authentifiée ou non), élément lu.
    \item Vérification~: confirmé, faux positif, non vérifiable.
    \item Vecteur CVSS et score, avec la mention B, BT, BE ou BTE.
    \item EPSS du jour~; inscription au catalogue KEV et date.
    \item Classe de priorité, échéance, responsable.
    \item Correction~: voie choisie, date, preuve de vérification.
\end{itemize}

\noindent Chapitres~\ref{chap:vulnerabilite} à~\ref{chap:gerer}.

\section*{M7. Demande d'exception}

\begin{itemize}
    \item Faille, classe, actif.
    \item Motif, avec sa preuve.
    \item Mesure compensatoire, sa date de mise en place et sa
          vérification.
    \item Durée et date de revue.
    \item Propriétaire du risque et décideur, conformes à la politique.
\end{itemize}

\noindent Chapitre~\ref{chap:gerer}.

\section*{M8. Plan d'un rapport d'expertise}

Mission (questions mot pour mot)~; pièces reçues~; opérations~;
constatations~; discussion~; limites~; réponses aux questions, cotées~;
annexes (Chapitre~\ref{chap:rapport}).

\section*{M9. Tableau des propositions cotées}

\begin{tableaucours}{cL{6.6cm}ccc}{%
    \tableauHeader{N°} &
    \tableauHeader{Proposition} &
    \tableauHeader{Date 1} &
    \tableauHeader{Date 2} &
    \tableauHeader{Date 3}
}
P1 & & & & \\
\end{tableaucours}

\noindent Avec, à chaque date, $\mu$, $U$ et $H$
(Chapitre~\ref{chap:auditologie}).
